Let $\triangle $ be the $n$-dimensional Laplacian, $\triangle
=\sum_{k=1}^n(\partial/\partial x_k)^2$. The fundamental solution $u$
for the equation $\triangle u=\delta$ can be found using the Fourier
transform. However, in this problem you are to fund the solution using
the following steps. By symmetry arguments, we can assume that $u$ is
a function $u(r)$ where $r=(x_1^2+\cdots+x_n^2)^{1/2}$. Rewrite the
equation $\triangle u=0$ away from $0$ as an ordinary differential
equation in the variable $r$. This is actually an Euler equation, as
studied in elementary differential equations. A solution can be
discovered by trying $u=r^s$. Explain how the constants appearing in
the solution can be evaluated, and evaluate them if you can. Explan
how the function $u$ can be used to solve the inhomogeneous (``Poisson'')
equation $\triangle u=f$.
